\chapter{Modèles de marge}\label{ch:rc:margin-models}

Au cours de mars 2020, la marge initiale exigée par les chambres de compensation du monde entier a augmenté d'environ 300
milliards de dollars, et les membres ont déposé 115 milliards de dollars de collatéral excédentaire en plus, soit environ 40~\%
plus que leur moyenne de février, appelée durant les semaines où chaque établissement avait besoin de liquidités pour autre chose
ailleurs. La hausse n'était pas une défaillance des modèles~: elle était ce que font les modèles qui suivent la volatilité lorsque
la volatilité bondit. La marge bilatérale calculée par le modèle standard du secteur n'a presque pas bougé, par construction.
La différence entre les deux tient à un choix en matière de procyclicité, et elle détermine qui doit trouver
des liquidités en période de crise. Le livre~1 a décrit la marge et la compensation~; ce chapitre construit les
modèles~: \omterm{def:rc:market-risk-measures:hs}{simulation historique} et \omterm{def:rc:market-risk-measures:fhs}{simulation historique filtrée} dans les chambres de compensation, les approches fondées sur les sensibilités et
standardisées pour les transactions bilatérales, les outils qui amortissent la procyclicité, et les tests
qui vérifient qu'un modèle de marge couvre ce qu'il doit couvrir.

\section{Ce que doit faire un modèle de marge initiale}

\begin{definition}[Modèle de marge initiale]\label{def:rc:margin-models:model}
Un \emph{modèle de marge initiale}\index{modèle de marge initiale} fixe le collatéral qu'une partie doit déposer
face à la perte que subirait sa contrepartie en liquidant son portefeuille après un défaut, sur
la période de marge en risque, avec un niveau de confiance élevé~: dans les chambres de compensation, au moins 99~\% unilatéral, pour
les transactions bilatérales, 99~\% sur dix jours.
\end{definition}

Un modèle de marge doit être sensible au risque (de sorte que la marge suive le portefeuille), stable (de sorte que
la marge ne bondisse pas à chaque mouvement de la volatilité), transparent et reproductible (afin que les membres
puissent anticiper leurs appels de marge), et rapide. Les deux premières exigences s'opposent l'une à l'autre.

\begin{example}[Le portefeuille de swaps d'un membre]\label{ex:rc:margin-models:book}
Les swaps compensés en USD d'un membre présentent des expositions illustratives de $+60$, $-40$, $-150$ et $-30$ milliers
de dollars par point de base à deux, cinq, dix et trente ans (le P\&L pour une hausse d'un point de base), soit un
net de $-160$ milliers de dollars. En prenant les taux au pair quotidiens du Trésor américain comme substituts des taux de swap, une \omterm{def:rc:market-risk-measures:hs}{simulation historique} de cinq jours
à 99~\% sur un an d'historique donne une marge initiale de 4.19~millions de dollars en moyenne en
février 2020.
\end{example}

\section{Modèles des chambres de compensation~: scénarios, simulation historique et simulation historique filtrée}

Les chambres de compensation calculent la marge par \omterm{def:rc:market-risk-measures:hs}{simulation historique} sur une fenêtre rétrospective, par grilles de scénarios (les
intervalles de balayage du livre~1), ou par \omterm{def:rc:market-risk-measures:fhs}{simulation historique filtrée}, qui remet à l'échelle les mouvements passés
à la volatilité courante (chapitre~21).

\omcode{code/firm/initmargin/firm_initmargin.py}{23}{46}{Mouvements pluriquotidiens qui se chevauchent, marge historique et filtrage à la volatilité du jour.}\label{lst:rc:margin-models:hs}

\begin{example}[Mars 2020 rejoué]\label{ex:rc:margin-models:replay}
En recalculant la marge du membre chaque jour de janvier à juin 2020, le modèle historique sur un an
passe de 4.19~millions de dollars (moyenne de février) à un pic de 5.25~millions de dollars le 18 mars, soit 25~\% de plus. Le
modèle filtré, avec un facteur de décroissance EWMA de 0.97, passe de 4.21~millions à 10.30~millions de dollars le 23 mars, soit 2.45
fois son niveau de février~: 6.09~millions de dollars de marge nouvelle en environ trois semaines
(la \cref{fig:rc:margin-models:replay}). Le modèle historique est lent, car une fenêtre d'un an dilue un
petit nombre de jours volatils~; le modèle filtré réagit tout de suite.
\end{example}

\begin{omfigure}
\begin{tikzpicture}
\begin{axis}[width=11.5cm,height=5.6cm,
  xlabel={jours depuis le 2 janvier 2020},ylabel={marge initiale (millions de dollars)},
  tick label style={font=\small},label style={font=\small},
  xmin=0,xmax=180,ymin=0,ymax=12,grid=major,grid style={gray!25},
  legend columns=2,legend cell align=left,legend style={font=\footnotesize,at={(0.5,-0.25)},anchor=north,draw=none,fill=none,/tikz/every even column/.append style={column sep=8pt}},
  table/col sep=comma]
\addplot[omExo,very thick] table[x=day,y=hs]{figdata/rates-credit-risk/25-margin-models/im2020.csv};
\addplot[omThm,very thick] table[x=day,y=fhs]{figdata/rates-credit-risk/25-margin-models/im2020.csv};
\addplot[omDef,thick,dashed] table[x=day,y=buffer]{figdata/rates-credit-risk/25-margin-models/im2020.csv};
\addplot[omMeth,thick,dotted] table[x=day,y=stressed]{figdata/rates-credit-risk/25-margin-models/im2020.csv};
\legend{historique (1 an),filtré (EWMA 0.97),filtré avec tampon de 25~\%,filtré avec pondération stressée de 25~\%}
\end{axis}
\end{tikzpicture}

\omcaption{Marge initiale du portefeuille de swaps du membre sur le premier semestre 2020 selon quatre
modèles. Le filtrage sur la volatilité courante plus que double la marge en mars~; un tampon ou une
pondération en période stressée relève la marge déjà détenue avant le stress et réduit ainsi la hausse, sans
modifier le pic. Données~: taux au pair du Trésor américain~; le tutoriel du chapitre.}\label{fig:rc:margin-models:replay}
\end{omfigure}

\begin{dated}{2026-09}{La marge en mars 2020}\label{dat:rc:margin-models:march2020}
Le Comité de Bâle, le CPMI et l'IOSCO ont indiqué en septembre 2022 que la marge initiale exigée dans
l'ensemble des chambres de compensation avait progressé d'environ 300 milliards de dollars en mars 2020, avec 115 milliards de dollars
supplémentaires de collatéral excédentaire, soit environ 40~\% au-dessus de la moyenne de février~; la marge bilatérale calculée selon le \omterm{def:rc:reduced-form-credit:isda}{modèle standard de l'ISDA}, utilisé
par la plupart des banques au titre des règles de marge, est restée relativement stable. Environ la moitié des chambres de compensation
interrogées ne disposaient d'aucun dispositif formel de lutte contre la procyclicité~; les seuils de tolérance déclarés pour les hausses de marge s'échelonnaient
de 25~\% à 80~\% sur cinq jours.
\end{dated}

\section{Le modèle bilatéral standard}

\begin{definition}[Modèle standard de marge initiale]\label{def:rc:margin-models:simm}
Un \emph{modèle standard de marge initiale}\index{modèle standard de marge initiale} calcule la marge bilatérale
à partir de sensibilités prescrites (delta, vega, courbure) par classe de risque, pondérées par des poids de risque
calibrés, puis agrégées selon des corrélations prescrites, $\mathrm{IM} = \sqrt{WS^\top R\,WS}$ au sein de chaque classe~: la
même structure que la \omterm{def:rc:regulatory-capital-for-trading-books:sbm}{méthode fondée sur les sensibilités} du FRTB (chapitre~23). Le modèle du secteur est le SIMM de l'ISDA,
dont les paramètres sont recalibrés périodiquement~; le chapitre utilise des paramètres illustratifs.
\end{definition}

\begin{definition}[Marge initiale par barème, seuil de marge initiale]\label{def:rc:margin-models:schedule}
La \emph{marge initiale par barème}\index{marge initiale par barème} applique des taux fixes
au notionnel selon la classe d'actifs et l'échéance (pour les taux d'intérêt, 1~\%, 2~\% et 4~\% pour des échéances allant jusqu'à deux ans, de deux
à cinq ans, et plus de cinq ans~; 6~\% pour le change~; 15~\% pour les actions), nettée selon $0.4\times\text{brut}+0.6\times\mathrm{NGR}\times\text{brut}$
avec NGR le ratio du coût de remplacement net sur le coût brut. Le \emph{seuil de
marge initiale}\index{seuil de marge initiale}, au plus 50~millions d'euros par groupe, est la marge qu'une partie
n'a pas besoin de collecter.
\end{definition}

\begin{example}[Trois chiffres bilatéraux]\label{ex:rc:margin-models:bilateral}
Aujourd'hui, pour le même portefeuille négocié de gré à gré~: une marge historique sur dix jours à 99~\% de 4.15~millions de dollars~; une
marge fondée sur les sensibilités de 8.98~millions de dollars, avec des poids de risque illustratifs sur dix jours de 60, 60, 55 et 50
points de base~; et une marge par barème de 13.19~millions de dollars brute, 9.23~millions de dollars nette pour un NGR de 0.5
(la \cref{fig:rc:margin-models:approaches}). La marge par barème ignore les effets de compensation entre les échéances~; le
modèle fondé sur les sensibilités ne les prend en compte que par ses corrélations.
\end{example}

\begin{omfigure}
\begin{tikzpicture}
\begin{axis}[width=11.5cm,height=5cm,xbar,bar width=12pt,
  symbolic y coords={schedule gross,schedule (NGR 0.5),sensitivity-based,HS 10-day},ytick=data,yticklabels={{barème brut},{barème (NGR 0.5)},{par sensibilités},{HS 10 jours}},
  xlabel={marge initiale (millions de dollars)},
  tick label style={font=\small},label style={font=\small},
  xmin=0,xmax=15,grid=major,grid style={gray!25},
  nodes near coords,every node near coord/.append style={font=\footnotesize,/pgf/number format/fixed,/pgf/number format/precision=2},
  table/col sep=comma]
\addplot[fill=omDef!60,draw=omDef] table[y=approach,x=im]{figdata/rates-credit-risk/25-margin-models/approaches.csv};
\end{axis}
\end{tikzpicture}

\omcaption{Marge initiale bilatérale du portefeuille du membre aujourd'hui selon un modèle historique à dix jours,
le modèle fondé sur les sensibilités avec des paramètres illustratifs et le barème réglementaire. Les méthodes plus simples
exigent davantage. Données~: taux au pair du Trésor américain~; le tutoriel du chapitre.}\label{fig:rc:margin-models:approaches}
\end{omfigure}

\begin{dated}{2026-09}{Règles de marge pour les dérivés non compensés}\label{dat:rc:margin-models:umr}
Dans le cadre du Comité de Bâle--OICV (tel que révisé en avril 2020), la marge initiale sur les dérivés
non compensés par une chambre de compensation couvre un mouvement unilatéral de 99~\% sur dix jours~; depuis le premier septembre 2022,
de façon permanente, elle s'applique entre groupes dont le notionnel agrégé moyen dépasse 8~milliards d'euros~;
le \omterm{def:rc:margin-models:schedule}{seuil de marge initiale} atteint jusqu'à 50~millions d'euros et le \omterm{def:rc:counterparty-exposure:csa}{montant minimal de transfert} jusqu'à
500\,000~euros. Le SIMM de l'ISDA en vigueur à compter du 11 juillet 2026 est la version 2.8, calibrée à décembre 2025, un
document public.
\end{dated}

\section{Procyclicité et ses outils}

\begin{definition}[Outil anti-procyclicité]\label{def:rc:margin-models:apc}
Un \emph{outil anti-procyclicité}\index{outil anti-procyclicité} (anti-procyclicality tool) limite la hausse de la marge en période de tensions
en constituant davantage en marché calme~: un coussin au-dessus de la marge du modèle, libérable en période de tensions, une
pondération sur une période de tensions dans la calibration, une longue fenêtre d'observation, ou un plancher sur la volatilité.
Les règles européennes pour les contreparties centrales exigent au moins l'un des éléments suivants~: un coussin d'au moins 25~\% de la
marge calculée, une pondération d'au moins 25~\% sur les observations de tensions, ou une marge non inférieure à
celle qu'une fenêtre d'observation de dix ans donnerait.
\end{definition}

\begin{definition}[Majoration pour concentration]\label{def:rc:margin-models:conc}
Une \emph{majoration pour concentration}\index{majoration pour concentration} (concentration add-on) impose une marge supplémentaire sur les positions importantes
par rapport à la capacité du marché à les absorber, car les liquider prendrait plus de temps que la
période de marge en risque et ferait bouger les prix.
\end{definition}

\omcode{code/firm/initmargin/firm_initmargin.py}{73}{90}{Deux outils anti-procyclicité~: un coussin libéré en période de tensions puis reconstitué, et une pondération de la période de tensions.}\label{lst:rc:margin-models:apc}

\begin{example}[Les outils de mars 2020]\label{ex:rc:margin-models:apc}
Avec un coussin de 25~\% sur le modèle filtré, libéré lorsque la volatilité EWMA du portefeuille dépasse deux fois sa
moyenne sur un an, le membre détenait 5.26~millions en février et 10.30~millions au pic~: soit 5.04~millions
de marge nouvelle au lieu de 6.09. Avec une pondération de 25~\% sur la marge de l'année la plus volatile avant 2020
(6.55~millions), il détenait 4.79~millions et devait en trouver 5.51~millions. Les outils déplacent la marge de la période de tensions
vers la période calme~; ils n'abaissent pas le pic.
\end{example}

\section{Backtest de la marge}

Un modèle de marge fait l'objet d'un backtest comme un modèle de VaR~: on compte les jours où la perte réalisée sur
la période de marge en risque a dépassé la marge détenue la veille. Les chambres de compensation publient leur couverture
et examinent les dépassements.

\begin{example}[Couverture depuis 2017]\label{ex:rc:margin-models:backtest}
Sur 2\,177 jours depuis janvier 2017, la perte à cinq jours du portefeuille du membre a dépassé la marge
historique sur un an 35 fois et la marge filtrée 37 fois, contre 22 attendus à 99~\%~: les deux modèles
couvrent moins qu'ils ne visent, surtout lors des épisodes de 2020 et de 2022, et les fenêtres qui se
chevauchent font se regrouper les dépassements.
\end{example}

\section{Tutoriel~: la marge à travers une crise}\label{tut:rc:margin-models:tutorial}

\begin{tutorial}
\textbf{Objectif.} Calculer la marge initiale du membre par \omterm{def:rc:market-risk-measures:hs}{simulation historique} et filtrée, rejouer le
premier semestre 2020 avec et sans \omterm{def:rc:margin-models:apc}{outils anti-procyclicité}, et comparer les approches bilatérales.
\textbf{État final~:} les nombres des
\cref{ex:rc:margin-models:replay,ex:rc:margin-models:bilateral,ex:rc:margin-models:apc,ex:rc:margin-models:backtest}
et les deux graphiques.
\begin{tutsteps}
\item \textbf{Données}~: \texttt{data()} lit les taux au pair~; \texttt{pnl()} applique les expositions.
\item \textbf{Rejeu}~: \texttt{replay\_2020()}.
\item \textbf{Bilatéral}~: \texttt{hs10\_today()}, \texttt{simm\_today()}, \texttt{schedule\_today()}.
\item \textbf{Backtest}~: \texttt{backtest\_since()}~; \texttt{fig\_rc\_margin.py} écrit les graphiques.
\end{tutsteps}
\textbf{Pour aller plus loin.} Utiliser une fenêtre d'observation de dix ans lorsque l'historique le permet~; ajouter un plancher de volatilité
à la moyenne décennale et comparer la hausse de mars.
\end{tutorial}

\section{Construction~: le moteur de marge initiale}\label{bld:rc:margin-models:initmargin}

\begin{build}
\bfield{Rôle} Les prévisions de marge de la firme~: ce que les chambres de compensation et les contreparties
bilatérales appelleront, pour la planification de liquidité, la MVA (chapitre~19) et les \omterm{def:rc:stress-testing-and-scenarios:stress}{tests de résistance} (chapitre~22).
\bfield{Interface} \texttt{overlapping}, \texttt{hs\_im}, \texttt{ewma\_vol}, \texttt{fhs\_moves}~;
\texttt{simm\_like(sens, tenors, rw)}~; \texttt{schedule\_rate}, \texttt{schedule\_im(trades, net,
gross)}~; \texttt{apc\_buffer}, \texttt{apc\_stressed\_weight}~; \texttt{backtest}, \texttt{peak\_increase}.
\bfield{Règles} Pertes positives~; mouvements sur plusieurs jours qui se chevauchent~; paramètres de sensibilité illustratifs~;
le barème et sa compensation NGR comme dans le cadre du Comité de Bâle--OICV.
\bfield{Tests d'acceptation} \texttt{code/firm/initmargin/tests/}~: sommes qui se chevauchent~; marge historique
d'un échantillon connu~; taux du barème et compensation~; marges de sensibilité à un seul facteur et compensatoires~; comportement du coussin
et de la pondération en période de tensions~; compte du backtest.
\bfield{Pour aller plus loin} Marge de réévaluation complète pour les options~; \omterm{def:rc:margin-models:conc}{majorations pour concentration} et pour liquidité~; le SIMM
sous licence avec ses fichiers de calibration~; marge intrajournalière.
\end{build}

\begin{omsources}
\item BCBS--CPMI--IOSCO, \emph{Review of margining practices}, septembre 2022.
\item BCBS--IOSCO, \emph{Margin requirements for non-centrally cleared derivatives}, avril 2020.
\item CPMI--IOSCO, \emph{Principles for financial market infrastructures}, avril 2012.
\end{omsources}

\section{Exercices}

\begin{exercise}[$\star$]\label{exo:rc:margin-models:1}
Calculer la marge du barème d'un swap de 100~millions de dollars américains à dix ans et d'un swap de
100~millions de dollars américains à un an, bruts.
\end{exercise}

\begin{exercise}[$\star$]\label{exo:rc:margin-models:2}
Pourquoi une marge historique sur un an réagit-elle lentement à un pic de volatilité~?
\end{exercise}

\begin{exercise}[$\star$]\label{exo:rc:margin-models:3}
Le coût de remplacement net d'une contrepartie est de 30, contre un brut de 120. De quelle part la compensation
réduit-elle sa marge du barème~?
\end{exercise}

\begin{exercise}[$\star\star$]\label{exo:rc:margin-models:4}
Pourquoi la marge SIMM bilatérale est-elle restée stable en mars 2020 alors que la marge de la chambre de compensation a augmenté~?
\end{exercise}

\begin{exercise}[$\star\star$]\label{exo:rc:margin-models:5}
Pourquoi les \omterm{def:rc:margin-models:apc}{outils anti-procyclicité} n'abaissent-ils pas le pic de marge dans le rejeu~?
\end{exercise}

\begin{exercise}[$\star\star$]\label{exo:rc:margin-models:6}
Contre quoi une \omterm{def:rc:margin-models:conc}{majoration pour concentration} protège-t-elle, que ne voit pas un modèle fondé sur la volatilité~?
\end{exercise}

\begin{exercise}[$\star\star\star$]\label{exo:rc:margin-models:7}
\emph{Programmation.} Calculer le nombre d'exceptions attendu et la p-valeur de Kupiec pour 35 exceptions en
2\,177 jours à 1~\%, en traitant les jours comme indépendants, puis expliquer pourquoi cette hypothèse est fausse.
\end{exercise}

\begin{exercise}[$\star\star\star$]\label{exo:rc:margin-models:8}
\emph{Trouver la faille.} «~Notre modèle de marge est bien calibré~: sa marge de février a couvert 99~\% des
mouvements de février.~»
\end{exercise}

\section{Problème~: les appels de marge de mars 2020}

\begin{problem}[{Problème du week-end --- trouver le cash}]\label{pb:rc:margin-models:1}
Le membre de ce chapitre compense son livre de swaps auprès d'une chambre de compensation
qui applique le modèle filtré. Sa trésorerie détient de la liquidité pour les appels de marge.

\medskip\noindent\textbf{Partie I --- Les appels.}
\begin{enumerate}
\item Donner la marge de février ainsi que le pic, avec sa date.
\item Donner le cash à trouver au cours de mars.
\item Qu'est-ce qui a porté la hausse~: les positions ou la volatilité~?
\item Quelle part de la hausse le modèle historique sur un an aurait-il appelée~?
\item Pourquoi la chambre de compensation pourrait-elle préférer le modèle filtré malgré tout~?
\end{enumerate}

\medskip\noindent\textbf{Partie II --- Les outils.}
\begin{enumerate}[resume]
\item Donner la marge de février et la hausse avec un coussin de 25~\%.
\item Les donner avec une pondération de 25~\% pour la période de stress.
\item Que coûte chaque outil à la contrepartie les années calmes~?
\item Quand faut-il reconstituer le coussin, et pourquoi lentement~?
\item Quelles tolérances déclarées pour les hausses sur cinq jours la revue de 2022 a-t-elle relevées~?
\end{enumerate}

\medskip\noindent\textbf{Partie III --- Bilatéral.}
\begin{enumerate}[resume]
\item Donner la marge bilatérale de la contrepartie selon les trois approches.
\item Pourquoi une contrepartie préférerait-elle compenser~?
\item Quel est l'effet du seuil de 50~millions d'euros sur ce livre~?
\item Comment la marge de variation interagit-elle avec la marge initiale en période de crise~?
\item Quel modèle faudrait-il utiliser pour prévoir les besoins de liquidité~?
\end{enumerate}

\medskip\noindent\textbf{Partie IV --- Jugement.}
\begin{enumerate}[resume]
\item La marge procyclique est-elle un défaut ou une caractéristique~?
\item Qui supporte le coût de l'anti-procyclicité~?
\item Comment dimensionner le coussin de liquidité destiné à la marge~?
\item Énoncer le \emph{résultat nommé}~: la hausse de la marge initiale avec et sans coussin de 25~\% et une pondération de stress,
ainsi que le cash que la contrepartie a dû trouver.
\item En une phrase~: qu'arbitre un modèle de marge~?
\end{enumerate}
\end{problem}

\section{Questions d'entretien}

\begin{interviewq}[$\star$ \iqroles{risque, trader}]\label{iq:rc:margin-models:1}
À quoi sert la marge initiale, et comment une chambre de compensation la calcule-t-elle~?
\end{interviewq}

\begin{interviewq}[$\star\star$ \iqroles{risque}]\label{iq:rc:margin-models:2}
Expliquez la procyclicité de la marge et citez trois outils pour la combattre.
\end{interviewq}

\begin{interviewq}[$\star\star$ \iqroles{développeur, risque}]\label{iq:rc:margin-models:3}
Comment fonctionne un modèle de marge bilatérale fondé sur les sensibilités~?
\end{interviewq}

\begin{interviewq}[$\star\star$ \iqroles{trader, banque}]\label{iq:rc:margin-models:4}
Pourquoi deux chambres de compensation peuvent-elles facturer une marge différente pour le même swap~?
\end{interviewq}

\begin{interviewq}[$\star\star\star$ \iqroles{chercheur, risque}]\label{iq:rc:margin-models:5}
Comment backtesteriez-vous un modèle de marge, et qu'est-ce qui le rend difficile~?
\end{interviewq}

\begin{interviewq}[$\star\star\star$ \iqroles{banque}]\label{iq:rc:margin-models:6}
Qu'a appris mars 2020 sur la marge et la liquidité~?
\end{interviewq}
